\section{Embedding context-free session types}
\label{sec:cfst-vs-algst}

% PT: this should be integrated in the related work paragraph to avoid repetition

% In this section we informally compare the metatheory
% of algebraic session types with that of context-free session types. Although
% these theories are closely related expressivity-wise, the abstraction mechanisms
% underlying algebraic and context-free session types are quite different:
% algebraic session types parameterize over behavior, whereas context-free session
% types abstract over the exchanged messages. Consider the example of the
% \lstinline{Stream} protocol presented in~\cref*{sec:param-prot}:

% \begin{lstlisting}
%     protocol Stream (a:P) = Next a (Stream a)
% \end{lstlisting}

% The direction of the communication is not settled until the parameter is
% instantiated and the protocol is materialized: a stream of integer values is
% governed by type $\TOut{\mathsf{Stream\,Int}}{\TEndT}$, while a process
% that receives integers \emph{ad aeternum} may be typed by
% $\TOut{\mathsf{Stream\,-Int}}{\TEndT}$. If we were to define a context-free
% counterpart for \lstinline|Stream|, we could think of types
% $\forall \alpha\colon\kindm. \mu x\colon \kinds. !\alpha; x$ or
% $\forall \alpha\colon\kindm. \mu x\colon \kinds. ?\alpha; x$ but both commit to
% a concrete direction. Protocol \lstinline{Stream} abstracts the direction
% through parameter \lstinline{a}. One could argue that these could be achieve
% with higher-order types, in the lines of
% \citet{DBLP:journals/corr/abs-2203-12877}, but, to the best of our knowledge,
% handling fully polymorphic sessions (read, assign kind $\kinds$ to polymorphic
% types) is still an open problem. For these reasons, we claim that algebraic
% session types are more expressive than context-free session types.

This section presents an informal investigation of two questions that relate \algst and
context-free session types (CFST
\cite{DBLP:journals/corr/abs-2106-06658}): can we map any \algst
program into an equivalent CFST program? And: can we map any CFST program to an
equivalent \algst program?

This investigation is necessarily informal because the features of
\algst and CFST do not line up. \algst is a minimalist mostly linear
version of System~F with products, sessions, and algebraic
protocols. CFST is also based on System~F, but has a richer kind
structure that enables unrestricted as well as linear features,
variant and record types, and the full slate of context-free session
primitives.

To enable a meaningful comparison,  we focus on the linear fragment of
CFST~\cite{DBLP:journals/corr/abs-2106-06658}
using the notation of \citet{DBLP:journals/corr/abs-2203-12877},
restricted to pairs and ignoring De Bruijn indices.
In \algst, we only consider types that are not polymorphic in
protocols because CFST has no equivalent. We add linear algebraic
datatypes to cope with variant types in CFST.

% We decided to 
% present the theory of algebraic session types in a simplified form: all types are linear, 
% we do not have records and we omitted the definition of $\datak$type definitions in 
% the formalization of our system (despite the obvious similarities with $\protocolk$ definitions). 
% For that reason, this comparison should 
% be considered in a loose, rather than in a strict, sense. In the context-free side, we 
% consider pairs instead of records as well.

% and ignore De Bruijn 
%indices (as used by Costa et al.~\cite{DBLP:journals/corr/abs-2203-12877}) because they would 
%unnecessarily burden the analysis.

%\begin{figure}[t!]
  \declrel{AlgST${}_0$ to CFST: Translation of algebraic to context-free session types}{$\algsttocfst{\typec T}=T$}
 
    \begin{mathpar}
    \algsttocfst{\TUnit} = \syntax{unit} \quad
    \algsttocfst{\TFun[]TU} = \algsttocfst {\typec{T}} \rightarrow \algsttocfst {\typec{U}} \quad
    \algsttocfst{\TPair TU} = \algsttocfst {\typec{T}} \times \algsttocfst {\typec{U}}\quad
    \algsttocfst{\TForall{\alpha}{\kindt} T} = \forall \alpha. \algsttocfst {\typec{T}}\quad
	\algsttocfst{\typec{\alpha}} = \alpha\\
	%% sessions
	\algsttocfst{\TIn{T}{U}}= ?\algsttocfst{\typec{T}}; \algsttocfst{\typec U} \quad 
	\algsttocfst{\TOut{T}{U}}= !\algsttocfst{\typec T}; \algsttocfst{\typec U} \quad 
	\algsttocfst{\TEndT} = \syntax{skip} \quad 
	\algsttocfst{\TEndW} = \syntax{skip}\quad
	\algsttocfst{\polarizedparam{-}{T}} = \syntax{dual}\, \algsttocfst{\typec T} \\ 
	\algsttocfst{\TPApp{P}{\overline T}} = \oplus\{ C_i\colon \algsttocfst{\tosession{T_i^1}{\TEndT}} ; \ldots;\algsttocfst{\tosession{T_i^{n_i}}{\TEndT}}\}_{i\in I}  \text{ when }
	\protocolk\ \typec{\proto\ \overline{\alpha}}~\blk{=}~\typec{\{{\termc{\constr_i}\ T_i^1\ldots T_i^{n_i}}\}_{i\in I} }
    \end{mathpar}
\caption{Translation of algebraic session types to context-free session types}
\label{fig:algst-to-cfst}
\end{figure}


\paragraph{Non-parameterized algebraic sessions into (linear) context-free sessions} 
The translation from 
algebraic session types that do not parameterize over protocol types
(AlgST${}_0$) to context-free session types 
(CFST)  is  given by $\algsttocfst{\typec T}$, which is defined on
types in normal form, only:
\begin{mathpar}
    \algsttocfst{\TUnit} = \syntax{unit} \quad
    \algsttocfst{\TFun[]TU} = \algsttocfst {\typec{T}} \rightarrow \algsttocfst {\typec{U}} \quad
    \algsttocfst{\TPair TU} = \algsttocfst {\typec{T}} \times \algsttocfst {\typec{U}}\quad
    \algsttocfst{\TForall{\alpha}{\kindt} T} = \forall \alpha:\kindtl. \algsttocfst {\typec{T}}\\
	\algsttocfst{\typec{\alpha}} = \alpha\qquad
	%% sessions
	\algsttocfst{\TIn{T}{U}}= ?\algsttocfst{\typec{T}}; \algsttocfst{\typec U} \qquad
	\algsttocfst{\TOut{T}{U}}= !\algsttocfst{\typec T}; \algsttocfst{\typec U} \qquad 
	\algsttocfst{\TEndT} = \syntax{skip} \qquad 
	\algsttocfst{\TEndW} = \syntax{skip}\\
	\algsttocfst{\TIn{\TPApp{P}{\overline U}}{S}}= P_U; \algsttocfst{\typec S} 
	\text{ where } 
	P_U \doteq \&\{ C_i\colon
	\algsttocfst{\TFIn {\overline{T_i[\overline U/\overline\alpha]}}{\TEndT}} \}_{i\in I}
	\text{ for }
\protocolk\ \typec{\proto\ \overline{\alpha}}~\blk{=}~\typec{\{{\termc{\constr_i}\ \overline{T_i}}\}_{i\in I} } \\
	\algsttocfst{\TOut{\TPApp{P}{\overline U}}{S}}= P_U; \algsttocfst{\typec S} 
	\text{ where } 
	%% \algsttocfst{\polarizedparam{-}{T}} = \syntax{dual}\, \algsttocfst{\typec T} \\
	P_U \doteq \oplus\{ C_i\colon
        \algsttocfst{\TFOut {\overline{T_i[\overline U/\overline\alpha]}}{\TEndT}} \}_{i\in I}
        \text{ for }
	\protocolk\ \typec{\proto\ \overline{\alpha}}~\blk{=}~\typec{\{{\termc{\constr_i}\ \overline{T_i}}\}_{i\in I} }
    \end{mathpar}

Functional types and 
message exchanges are naturally mapped onto their counterparts. The end of a session is mapped 
to $\syntax{skip}$, whereas algebraic protocols are translated to
type names defined as internal choices or branches, depending on whether they appear 
in sending or receiving positions. The choice labels coincide with the protocol
constructors. 

We claim that if $\isEquiv{T}{U}$, then $\algsttocfst{\typec T}\simeq\algsttocfst{\typec U}$.

\begin{figure}[t!]
  \declrel{CFST to AlgST: Translation of the functional fragment}{$\cfsttoalgst{T\colon \mathsf{T^{lin}}}=\typec T \colon \kindt$}
 
    \begin{mathpar}
    \cfsttoalgst{\syntax{unit}} = \TUnit \and
    \cfsttoalgst{T \rightarrow U} = \TFun[]{\cfsttoalgst{T}}{\cfsttoalgst{U}} \and
    \cfsttoalgst{T \times U} = \TPair{\cfsttoalgst{T}}{\cfsttoalgst{U}}\\
%    \cfsttoalgst{\forall\alpha.T} = \TForall{\alpha}{\kindt}{\cfsttoalgst{T}}\quad
%	\cfsttoalgst{\alpha} = \typec{\alpha} \\
	\cfsttoalgst{\langle \ell \colon T_\ell\rangle_{\ell\in L}} = \typec{X_L }\text{ where }\datak\ \typec{X_L}~\blk{=}
	~\typec{\{{\termc{\ell}\ \cfsttoalgst{T_\ell}}\}_{\ell\in L} }
\\
	\cfsttoalgst{X\colon \kappa} = \typec{X}
	\text{ for } X \doteq T \text{ and } \kappa \neq \mathsf{S^{lin}}
	\text{ where }
	\datak\ \typec{X}~\blk{=}~\typec{\termc{\mathsf{Unfold}_X}\ \cfsttoalgst{T}}\\
	\cfsttoalgst{T\colon \mathsf{S^{lin}}} = \TOut{X_T}{\TEndT} \text{ where }\protocolk\ \typec{X_T}~\blk{=}~\typec{\termc{X_T}\ \cfsttoprot T}\text{ and $X_T$ is a fresh name.}
%% old version, with polymorphic types
%		\datak\ \typec{X\ \overline{\alpha}}~\blk{=}~\typec{\termc{\mathsf{Unfold}_X}\ \cfsttoalgst{T}} \text{ and }\FVL{T} = \overline\alpha\\
%	\cfsttoalgst{\langle \ell_1 \colon T_1, \ldots, \ell_n \colon T_n\rangle} = \TPApp {X_{\langle\ell_1\cdots \ell_n\rangle}}{\overline{\alpha}} \text{ where }\datak\ \typec{X_{\langle\ell_1\cdots \ell_n\rangle}\ \overline{\alpha}}~\blk{=}~\typec{\termc{\ell_1}\ \cfsttoalgst{T_1} \mid \ldots \mid \termc{\ell_n}\ \cfsttoalgst{T_n}} 
%	\text{ and }\cup_{i=1}^n \FVL{T_i} = \overline\alpha
%\\
%	\cfsttoalgst{T\colon \mathsf{S^{lin}}} = \TOut{X_T}{\TEndT} \text{ where }\protocolk\ \typec{X_T}~\blk{=}~\typec{\termc{X_T}\ \cfsttoprot T}\text{ and $X_T$ is a fresh name.}
    \end{mathpar}
    
    \declrel{CFST to AlgST:Translation of the session typed fragment}{$\cfsttoprot{T\colon \mathsf{S^{lin}}}=\typec{\overline {T\colon \kindp}}$}
    
    \begin{mathpar}
    \cfsttoprot{\syntax{skip}} = \emptyseq \and \cfsttoprot{! T} = \cfsttoalgst{T} \and
    \cfsttoprot{? T} = \TMinus {\cfsttoalgst{T}}\and
    \cfsttoprot{T; U} = \cfsttoprot{T} \cfsttoprot{U}
    	\\
	\cfsttoprot{\oplus\{\ell\colon T_\ell\}_{\ell\in L}} = \typec{X_L}
        \text{ where } \protocolk\ \typec{X_L}~\blk{=} ~\typec{\{{\termc{\ell}\ \cfsttoprot{T_\ell}}\}_{\ell\in L}}
	\\
	\cfsttoprot{\&\{\ell\colon T_\ell\}_{\ell\in L}} = \TMinus
        {\typec{X_L}}
        \text{ where } \protocolk\ \typec{X_L}~\blk{=}
        ~\typec{\{{\termc{\ell}\ \cfsttoprot{\syntax{dual}\,T_\ell}}\}_{\ell\in L}}
    \\
	\cfsttoprot{X\colon \mathsf{S^{lin}}} = \typec{X}
	\text{ for } X \doteq T
	\text{ where }
	\protocolk\ \typec{X}~\blk{=}~\typec{\termc{\mathsf{Unfold}_X}\ \cfsttoprot{T}} 
	%
	%~\typec{\termc{\ell_1}\ \cfsttoprot{T_1} \mid \ldots \mid \termc{\ell_n}\ \cfsttoprot{T_n}} 
    \end{mathpar}
\caption{Translation of linear context-free session types to algebraic session types}
\label{fig:cfst-vs-algst-f}
\end{figure}


%%% Local Variables:
%%% mode: latex
%%% TeX-master: "main"
%%% End:


\paragraph{Linear context-free sessions embedded into algebraic sessions} 
In this direction, we omit polymorphism to obtain a lighter, more
perspicuous notation. Our comparative approach scales to the
polymorphic setting, but at the price of high notational overhead as
pointed out in \cref{sec:appendix-cfst-algst}. 

The translation of the \emph{linear} and \emph{monomorphic} fragment of CFST 
to AlgST is split between the functional and the session typed fragments,
which are mutually defined and presented in~\cref{fig:cfst-vs-algst-f}. The translation of a
context-free session type $T$ to the algebraic setting is given by $\cfsttoalgst T$. 
The functional operators have natural counterparts in the algebraic side.
When translated to the isorecursive setting, functional recursive types are represented 
by a datatype enclosed with an explicit unfold. Variant types are mapped to datatypes whose 
constructors are the respective labels. Session types are materialized by protocol definitions:
each session type originates a protocol emerging from the conversion of a sequential composition of 
types into a sequence of (algebraic) types with kind $\kindp$. 
This conversion is governed by function $\cfsttoprot{\cdot}$ that maps $\syntax{skip}$
to the empty sequence and uses polarities to enforce the direction of communication.
In the session typed fragment, recursive types and choices are represented by protocol definitions.
Type names are considered generative and we relax \algst to admit 
overloaded constructor names that can be reused across protocol definitions.

%\begin{table}
\begin{tabular}{| c | l |}
\hline
Equivalence rule & Isomorphism witness\\
$T\simeq U \colon \mathsf{S^{lin}}$&$\cfsttoalgst{T} \to
\TDual{\cfsttoalgst{U}} \to \TUnit$\\\hline
%% E-Skip
$\axiom{}{\isodecl{\delta} \mathsf{Skip}\simeq \mathsf{Skip}}$ & 
\begin{lstlisting}
delta t u = wait u 
        |> terminate t  
\end{lstlisting}
\\\hline
 %% E-IdL
$\inferrule{X\doteq T \\ T\,\mathsf{contr}\\\isodecl{\delta_1} T \simeq U}{\isodecl{\delta} X \simeq U}$ & 
\begin{lstlisting}
delta x u = let t = select UnfoldX x in 
        delta1 t u
\end{lstlisting}
 \\\hline
 %% E-IDSeqL
$\inferrule{X\doteq T \\ T\,\mathsf{contr}\\\isodecl{\delta_1} T;V \simeq U}{\isodecl{\delta} X;V \simeq U}$ & 
\begin{lstlisting}
delta x u = let t = select UnfoldX x in 
        delta1 t u
\end{lstlisting}
 \\\hline
 %% E-MsgSeq2
$\inferrule{\isodecl{\theta} T \simeq U\\\isodecl{\delta_1} V \simeq W}{\isodecl{\delta} !T; V \simeq !U;W}$ & 
\begin{lstlisting}
delta v w = let (u, w) = receive w in 
        let t = theta^-1 u 
        let v = send t v in 
        delta1 v w
\end{lstlisting}
 \\\hline
 %% E-MsgSeq1L
$\inferrule{\isodecl{\theta} T \simeq U\\\isodecl{\sigma} V \checkmark}{\isodecl{\delta} !T; V \simeq !U}$ & 
\begin{lstlisting}
delta t u = let (u', u) = receive u in 
        let t' = theta^-1 u'
        let v = send t' t in
        let v = sigma v in
        wait u 
        |> terminate v
\end{lstlisting}
 \\\hline
%% E-Msg
$\inferrule{\isodecl{\theta} T \simeq U }{\isodecl{\delta} !T \simeq !U}$ & 
\begin{lstlisting}
delta t u = let (u', u) = receive u in
        let t' = theta^-1 u' in
        let t = send t' t in
        wait u 
        |> terminate t
\end{lstlisting}
 \\\hline
%% E-Choice
$\inferrule{\isodecl{\delta_\ell} T_\ell \simeq U_\ell\\ (\forall \ell \in L)}{\isodecl{\delta} \oplus\{\ell\colon T_\ell\}_{\ell \in L} \simeq \oplus\{\ell\colon U_\ell\}_{\ell \in L}}$ & 
\begin{lstlisting}
delta t u = match u with {
          k u -> let t = select k t in 
                 deltak t u }
\end{lstlisting}
 \\\hline
%% E-ChoiceSeqL
$\inferrule{\isodecl{\delta_1} \oplus\{\ell\colon T_\ell;U\}_{\ell \in L} \simeq V}{\isodecl{\delta} \oplus\{\ell\colon T_\ell\}_{\ell \in L};U \simeq V}$ & 
\begin{lstlisting}
delta t u = delta1 t u
\end{lstlisting}
 \\\hline
 %% E-SkipSeqL
$\inferrule{\isodecl{\delta_1} T \simeq U}{\isodecl{\delta} \mathsf{Skip}; T \simeq U}$ & 
\begin{lstlisting}
delta t u = delta1 t u
\end{lstlisting}
 \\\hline
% %% E-IndexSeq1L
%$\inferrule{\isodecl{\sigma} T\checkmark}{\isodecl{\delta} n;T \simeq n}$ & 
%\begin{lstlisting}
%delta t u = ?
%\end{lstlisting}
% \\\hline
    %%
\end{tabular}
\caption{Processes that consume two dual session endpoints.}
\label{tab:iso-witness-sessions}
\end{table}

\begin{table}
\begin{tabular}{| c | l |}
\hline
Equivalence rule & Isomorphism witness\\
$T\simeq U \colon \kappa$ with $\kappa\neq \mathsf{S^{lin}}$&$\theta \colon \cfsttoalgst{T} \rightarrow \cfsttoalgst{U}$\\\hline
%% E-Unit
$\axiom{}{\isodecl{\theta} ()_{\mathsf{lin}}\simeq ()_{\mathsf{lin}}}$ & 
\begin{lstlisting}
theta = lambda x:Unit.x
\end{lstlisting}
\\\hline
%%% E-Index
%$\axiom{}{\isodecl{\theta} n}\simeq n}$ & 
%\begin{lstlisting}
%theta = lambda x:Unit.x
%\end{lstlisting}
%\\\hline
%% E-Arrow
$\inferrule{\isodecl{\theta_1} T \simeq U\\\isodecl{\theta_2} V \simeq W}{\isodecl{\theta} T\rightarrow V \simeq U \rightarrow W}$ & 
\begin{lstlisting}
theta t = lambda x:[|U|].theta2 $ t $ theta^-1 x
\end{lstlisting}
 \\\hline
%% E-Rcd
$\inferrule{\isodecl{\theta_1} T \simeq U\\\isodecl{\theta_2} V \simeq W}{\isodecl{\theta} T\otimes V \simeq U \otimes W}$&
\begin{lstlisting}
theta t = <theta1 $ fst t,theta2 $ snd t> 
\end{lstlisting}
 \\\hline
  %% E-IdL
$\inferrule{X\doteq T \\ T\,\mathsf{contr}\\\isodecl{\theta_1} T \simeq U}{\isodecl{\theta} X \simeq U}$ & 
\begin{lstlisting}
theta x = let t = select UnfoldX x in 
      let u = theta1 t in
      u
\end{lstlisting}
%  \\\hline
%  %% E-Quant
% $\inferrule{\isodecl{\theta_1} T \simeq U}{\isodecl{\theta} \forall T \simeq \forall U}$&
% \begin{lstlisting}
% theta t = Lambda alpha: T. theta1 $ t alpha
% \end{lstlisting}
 \\\hline
    %%
\end{tabular}
\caption{Transformation of non-session typed expressions.}
\label{tab:iso-cfst-to-algst-f}
\end{table}

\begin{table}
\begin{tabular}{| c | l |}
\hline
Is terminated predicate & Witness\\
$T \colon \mathsf{S^{lin}}$&$\sigma \colon \cfsttoalgst{T} \rightarrow \TEndT$\\\hline
%% E-Unit
$\axiom{}{\isodecl{\sigma} \mathsf{Skip}\checkmark}$ & 
\begin{lstlisting}
sigma t = t
\end{lstlisting}
\\\hline
%% E-Arrow
$\inferrule{\isodecl{\sigma_1} T \checkmark\\\isodecl{\sigma_2} U \checkmark}{\isodecl{\sigma} T;U \checkmark}$ & 
\begin{lstlisting}
sigma t = let t' = sigma1 t in
      sigma2 t'
\end{lstlisting}
 \\\hline
%% E-Rcd
$\inferrule{T\doteq U\\ \isodecl{\sigma_1} T \checkmark}{\isodecl{\sigma} X \checkmark}$&
\begin{lstlisting}
sigma x = let t = select UnfoldX x in
      sigma1 t
\end{lstlisting}
 \\\hline
    %%
\end{tabular}
\caption{Processes that consume the explicit unfolds created by the translation.}
\label{tab:iso-is-terminated}
\end{table}

%%% Local Variables:
%%% mode: latex
%%% TeX-master: "main"
%%% End:


It remains to define suitable adapters in the places where CFST relies
on type equivalence of equirecursive types. To this end we generate
isomorphisms as witnesses of equivalence proofs using the rule system
proposed by \citet{DBLP:journals/corr/abs-2203-12877}. 
The tricky part of constructing such an isomorphism is its extension
into session types. 
For simplicity, we consider just one side of the isomorphism
$\mathsf{iso} : \cfsttoalgst{T} \rightarrow \cfsttoalgst{U}$ 
where $T \simeq U$ as session types:
\begin{lstlisting}
iso t = let (u, u') = new [|U|] in fork $ delta t u' |> u
\end{lstlisting} %$
This definition relies on a process $\delta : \cfsttoalgst{T} \to
\TDual{\cfsttoalgst{U}} \to \TUnit$ that consumes two (dual) channel
ends and implements a forwarder that compensates for the differences
between $T$ and $U$. The witness processes $\delta$ are 
defined for each pair of equivalent session types, according 
to the rules. We sketch some witness 
processes in \cref{sec:appendix-cfst-algst} and provide a more detailed analysis of the isomorphism. 

In general, $\delta$ inserts the explicit unfolds imposed by the translation to the isorecursive setting.
We showcase our approach for the equivalence $Y \simeq U$ where type $Y$ is defined by $T$, $Y \doteq T$.
The definition of $\delta$ follows the equivalence rule of \citet{DBLP:journals/corr/abs-2203-12877}
and shown below.
%
A witness process $\delta$ for the embedding $\mathsf{iso} : \cfsttoalgst{Y} \rightarrow \cfsttoalgst{U}$  
builds on a process $\delta_1$ witnessing $T \simeq U$ and is defined by:

\begin{minipage}{0.6\textwidth}
\begin{lstlisting}
delta y u = match u with {
   XU u -> let y = select XY y in 
           let t = select UnfoldY y in
           delta1 t u }
\end{lstlisting}
\end{minipage}
\begin{minipage}{0.4\textwidth}
	$$\inferrule{Y\doteq T \\ T\,\mathsf{contr}\\\isodecl{\delta_1} T \simeq U}{\isodecl{\delta} Y \simeq U}$$
\end{minipage}

As a result of the definition of $\cfsttoalgst{Y\colon \mathsf{S^{lin}}}$, 
process $\delta$ \lstinline|match|es with $X_U$ on channel $u$ and \lstinline|select|s $X_Y$ on channel $y$, as prescribed by the
protocol definitions resulting from $\cfsttoalgst{\syntax{Y}}$ (see~\cref{fig:cfst-vs-algst-f}). Then, 
 \lstinline|select|s the explicit $\termc{\mathsf{Unfold}_Y}$ introduced by $\cfsttoprot{Y}$ and 
uses $\delta_1$ to consume the leftovers at $t$ and $u$.

%Other cases and a more thorough analysis can be found in \cref{sec:appendix-cfst-algst}. 
For the
other direction of the isomorphism, we could follow a similar approach for defining 
$\mathsf{iso}^{-1} : \cfsttoalgst{U} \rightarrow \cfsttoalgst{T}$ 
and then prove that $\mathsf{iso}$ and $\mathsf{iso}^{-1}$ compose giving the identity.
To extend the translation to the polymorphic fragment, we would need to consider a 
\emph{reader monad} to store a map from type variables to the auxiliary processes
that consume the respective channel ends.
 
% In the case of $!T; V \simeq !U;W$, we are provided not only with a witness $\delta_1$ to consume the remainder
% of channels $v$ and $w$, but also with a function $\theta \colon \cfsttoalgst{T} \rightarrow \cfsttoalgst{U}$
% that witnesses the isomorphism resulting from $T\simeq U$. The functions witnessing the isomorphism 
% between types in the functional fragment are presented in~\cref{tab:iso-cfst-to-algst-f}.
% Given $\delta_1$ and $\theta$, process $\delta$ \lstinline|receive|s $u$ on channel $w$, 
% transforms $u$ with $\theta^{-1}$ and sends the result through channel $v$. Process
% $\delta_1$ consumes the remainder of $v$ and $w$.

% To witness the isomorphism corresponding to $!T; V \simeq !U$ where $V$ is a 
% terminated session (see \citet{DBLP:journals/corr/abs-2203-12877} for further details), we are given
% a process $\sigma$ that consumes the explicit unfolds resulting from the isorecursive interpretation of $V$. The 
% \emph{consumer} processes $\sigma$ for terminated sessions are presented in \cref{tab:iso-is-terminated}.

% The other cases in \cref{tab:iso-cfst-to-algst-f,tab:iso-witness-sessions,tab:iso-is-terminated} follow
% a similar reasoning. The omitted cases for monomorphic types also fit the same approach. For the polymorphic 
% fragment, we would need to consider a \emph{reader monad} where one could store a map from type
% variables to the processes that consume the corresponding channels or transform the 
% corresponding values, depending on whether the type variables are of kind $\kinds$ or not. 
% In the axiom for polymorphic (non-session typed) variables $n\simeq n$, for instance,
% %\begin{mathpar}
% %\axiom{}{\isodecl{\theta} n}\simeq n}
% %\end{mathpar}
% we could think of $\theta$ as the identity function, but we wouldn't know which type to 
% assign to the bound variable: $\theta = \abse{x}{?}{x}$. Whereas to consume two channels 
% typed with a polymorphic session-typed variable, we would need to know how to consume both
% endpoints.
% %
% %$\inferrule{\isodecl{\sigma} T\checkmark}{\isodecl{\delta} n;T \simeq n}$
% %
% Processes to consume or transform values typed with polymorphic variables would need to 
% be provided to $\delta$. For the translation of types in the polymorphic fragment, we would need to collect the free
% variables occurrig in the type and use a parametrized $\datak$ or $\protocolk$ definition:
% \begin{mathpar}
% \cfsttoalgst{X\colon \kappa} = \typec{X\ \overline{\alpha}}
% 	\text{ for } X \doteq T \text{ and } \kappa \neq \mathsf{S^{lin}}
% 	\text{ where } \datak\ \typec{X\ \overline{\alpha}}~\blk{=}~\typec{\termc{\mathsf{Unfold}_X}\ \cfsttoalgst{T}} \text{ and }\FVL{T} = \overline\alpha
% 	\and 
% 	\cfsttoprot{X\colon \mathsf{S^{lin}}} = \typec{X\ \overline{\alpha}}
% 	\text{ for } X \doteq T
% 	\text{ where }
% 	\protocolk\ \typec{X\ \overline{\alpha}}~\blk{=}~\typec{\termc{\mathsf{Unfold}_X}\ \cfsttoprot{T}} 
% 	\text{ and }\FVL{T} = \overline\alpha
% \end{mathpar}
% %
% Although we believe that polymorphic types
% do not pose any problem to this analysis, we decided to keep the notation simple and
% focus on the monomorphic fragment.

% In this analysis we only provided witnesses for the "isomorphism" in one direction.
% For the other direction, we could follow a similar approach and then prove that 
% they are inverses of each other.

%%% Local Variables:
%%% mode: latex
%%% TeX-master: "main"
%%% End:
